% ============================================================
%  sec-1-introduction.tex  —  1.1 Introduction
% ============================================================

\section{Introduction}
\label{sec:introduction}

The rapid evolution of artificial intelligence has driven an unprecedented
increase in the complexity and scale of machine learning models \cite{brown2020gpt3}. Modern
applications — particularly large language models and multimodal systems —
demand significant computational resources \cite{kaplan2020scaling}, historically supplied by
NVIDIA's \acrshort{cuda}-based \acrshort{gpu} infrastructures. However,
the current AI hardware landscape is undergoing a structural transformation,
with specialised accelerators such as Neural Processing Units (\acrshort{npu}s)
emerging as viable alternatives optimised for efficiency, energy consumption,
and scalability \cite{iea2024electricity}.

This transition introduces substantial technical challenges: differences in
execution models, framework compatibility, model conversion processes, and
optimisation strategies make it increasingly difficult for organisations to
determine whether a given AI workload can be efficiently migrated — and at
what performance or cost implication.

This project addresses the problem by proposing an \term{intelligent
decision-support system} that analyses AI workloads and evaluates their
feasibility for migration across heterogeneous hardware environments.
Specifically, the system focuses on the transition from \acrshort{cuda}-based
\acrshort{gpu} environments to Huawei's Ascend \acrshort{npu} ecosystem,
providing structured insights into model compatibility, expected performance
variation, and infrastructure suitability.

The proposed platform introduces a modular architecture separating workload
analysis, hardware abstraction, migration evaluation, and recommendation
generation. Rather than relying solely on traditional benchmarking, the system
integrates benchmarking as a validation mechanism within a broader migration
intelligence framework, enabling decision support that goes beyond performance
measurement toward informed deployment choices.

The work is carried out within \CompanyName, leveraging the company's expertise
in AI infrastructure, cloud computing, and hardware acceleration technologies \cite{huawei2025annual}.
Huawei's Ascend ecosystem — together with its software stack including
\acrshort{cann} and MindSpore — provides the industrial foundation for
exploring alternative AI execution environments and future-oriented deployment
strategies.
